La adopción de servicios tecnológicos en el agro colombiano va más allá de la mera disponibilidad de aplicaciones y dispositivos; implica comprender las actitudes de los agricultores hacia el cambio y su capacidad económica. Este artículo analiza las barreras generacionales, económicas y culturales que enfrentan los productores rurales, mayoritariamente mayores de 50 años, y propone estrategias para una transformación digital inclusiva a través del Portal Agro Comercial del Huila. Se enfatiza en el diseño amigable, la capacitación presencial y la integración de funcionalidades offline para eliminar la brecha digital y promover el comercio justo, la trazabilidad y la sostenibilidad en la agricultura colombiana.